\scanpage{118}
% archive sheet number: 117
$\boxed{18}$ L'ens \deleted{\unsure{des}} $S$ engendre affinement $X$ fibre par fibre

[$\boxed{19}$ condition de connexité pour l'ens ordonné $\Psi_*(P)$, impliquant que $\Gamma_{\deleted{\unsure{P}}}$ opère transitivement sur $\mathrm{Rep}(P)$ \inserted{qui est un quotient de $\Gamma$ par suite (grâce à (18), (17))}, et que tt autom de $(X,P)$ \inserted{--- et $=$ tt autom combinatoire de $P$ ---} qui fixe un repère est l'identité]
% slanted
\margin{redondant ? Ou au contraire faut-il exiger en $=$ temps que $\Psi_*(P)/\Gamma_H\to\Psi_H(P)$ soit injectif $\forall H\subset\{0,1,2\}$, $H\ne\varnothing$ ?}

$\boxed{20}$ $\mathrm{Aut}(X,P)$ \uline{opère transitivement sur} $\mathrm{Rep}(P)$.

On trouve alors que le stabilisateur \inserted{de $\Gamma$} d'un repère $r$ est un sous-groupe \uline{invariant} $N$ de $\Gamma$, donc $\Psi_*(P)$ est un \uline{torseur} sous
\[
 \boxed{\Gamma/N=\unsure{\Gamma_P^\sharp|\Gamma'_P}}\tag*{(21)}
\]
et que
\[
 \boxed{\mathrm{Aut}(X,P)\simeq\mathrm{Aut}_{\mathrm{comb}}(P)}\tag*{(22)}
\]
est son commutant. Donc le choix d'un repère $r$ définit un isomorphisme
\[
 \deleted{\unsure{\Phi|\mathscr{C}}=}\Gamma_P\xrightarrow[\sim]{\varphi}C=\mathrm{Aut}(X,P)\tag*{(23)}
\]
par \inserted{pour $\gamma\in\Gamma$,} $\varphi(\gamma)$ est l'unique autom de $(X,P)$ tel que
\[
 \varphi(\gamma)(r)=\gamma.r_{\unsure{0}}\tag*{(24)}
\]

\uline{Question} \deleted{\unsure{Pour application}} $r\in\mathrm{Rep}(P)$ fixé, l'application
\[
 \Gamma\longrightarrow\mathrm{Rep}(P)\qquad \gamma\longmapsto\gamma.r\tag*{(25)}
\]
est surjective par (19), et \deleted{\unsure{induit}} par passage aux quotients des
\[
 N\backslash\Phi_H\deleted{/\unsure{N}}=\Gamma_P/\overline\Gamma_H=\Gamma/N.\Gamma_H\longrightarrow\Psi_H(P)\tag*{(26)}
\]
